% !TeX root = ../教材题库_全部条目_转写工作稿.tex
% Source: XJ-B1-C04-RV-P153-Y1216
% 第4章 / 章末复习 复习题四 / 教材 p.153、154 / PDF 第160、161页
% OCR未校对：公式与文字均可能有误；未完成数学TeX转写；禁止直接发给学生
% 原题及相邻原解截图（核对用）：../题目原图/XJ-B1-C04-RV-P153-Y1216-01.webp, ../题目原图/XJ-B1-C04-RV-P153-Y1216-02.webp
% 本文件只在转写完成后作为题干使用；OCR 条目尚不是可靠题干。
\begin{bankitem}{复习题四 第25题}{第4章 / 章末复习 复习题四 / 教材 p.153、154 / PDF 第160、161页}{XJ-B1-C04-RV-P153-Y1216}{OCR未校对：公式与文字均可能有误；未完成数学TeX转写；禁止直接发给学生}
\begin{ocrdraft}
[教材p.153 / PDF第160页]
\par
25．将温度探头放入一杯水中，随着时间变化记录水温数据，得到下表数据．
\par
时间/mm
\par
1
\par
2
\par
3
\par
4
\par
5
\par
温度/\ensuremath{{}^\circ\mathrm C}
\par
90．5
\par
82．5
\par
75．5
\par
69．0
\par
63．5
\par
(1)描点画出水温随时间变化的大致图象；
\par
（2）建立一个能基本反映水温3，（\ensuremath{{}^\circ\mathrm C}）随时间x(mrn)变化的函数模型，并借助
\par
计算机软件作出其图象，观察它与描点画出的图象的吻合程度；
\par
（3）分析所得的函数模型图象，估计经过多少分钟水温才会降到35\ensuremath{{}^\circ\mathrm C}左右．
\par
[教材p.154 / PDF第161页]
\par
，穿丆
\par
22
\par
0入
\par
0\ensuremath{\text{\textperthousand}}入
\end{ocrdraft}
\end{bankitem}
